\documentclass[conference]{IEEEtran}
\IEEEoverridecommandlockouts

\usepackage{fancyhdr}
\pagestyle{fancy}
\fancyhf{}

\renewcommand{\thepage}{\Roman{page}} % Capital Roman numerals

\fancyhead[L]{\thepage}                 % Page number on left
\fancyhead[C]{S. VERMA \& S. DHANUKA}        % Standard author format

% --------------------------------
% -----------------------------------

\usepackage{cite}
\usepackage{amsmath,amssymb,amsfonts}
\usepackage{algorithmic}
\usepackage{graphicx}
\usepackage{textcomp}
\usepackage{xcolor}
\usepackage{booktabs}
\usepackage{multirow}

\begin{document}

\title{Determinants of Household Participation in Securities Markets: An Empirical Analysis\\
\thanks{Based on the 2025 SEBI Household Survey on Securities Market Participation.}
}

\author{\IEEEauthorblockN{Shivansh Verma}
\IEEEauthorblockA{\textit{Dept of CS} \\
\textit{Ashoka University}\\
Sonipat, India \\
shivansh.verma\_ug2023@ashoka.edu.in}
\and
\IEEEauthorblockN{Sashwat Dhanuka}
\IEEEauthorblockA{\textit{Dept of CS} \\
\textit{Ashoka University}\\
Sonipat, India \\
sashwat.dhanuka\_ug2023@ashoka.edu.in}
}

\maketitle

\thispagestyle{empty} % No header on first page


\begin{abstract}
Studies of household participation in securities markets routinely attribute low retail investment rates to income and education constraints. We test this assumption using the 2025 SEBI Household Survey ($N = 109,430$), a nationally representative cross-section administered across four geographic zones in India. A survey-weighted logistic regression finds that product awareness (the count of financial instruments a household recognizes) is the single strongest predictor of participation (OR = 2.22), exceeding income (OR = 1.60) and education (OR = 1.56) even after demographic controls. This implies that the dominant barrier to entry is not poverty but unfamiliarity with the instrument landscape, and that awareness-focused outreach is a distinct policy lever from income-based redistribution. Among non-investors, K-means clustering on self-reported barriers reveals three archetypes with different intervention needs: Trust-Deficient (46.3\%), Knowledge-Gated (29.0\%), and Fear-Driven (24.8\%). Once inside the market, demographic predictors lose explanatory power; information source and financial self-concept become the primary axes of behavioral segmentation. Investors relying on social media finfluencers concentrate in speculative assets and short-term motives, while overconfident investors (high perceived familiarity, low objective knowledge) disproportionately hold Futures \& Options and cryptocurrency. Gradient Boosting Classifiers confirm these patterns predictively (participation AUC = 0.814, securities choice AUC = 0.820). The findings suggest that broadening market participation in India requires parallel interventions: expanding awareness to unlock entry, and improving information quality to correct post-entry behavior.
\end{abstract}


\section{Introduction}
The democratization of capital markets is essential for broad-based economic growth and household wealth accumulation. In emerging economies like India, retail participation in securities has witnessed a significant surge, catalyzed by digitization and the proliferation of zero-commission trading platforms. Despite this growth, participation remains highly heterogeneous.

Understanding the determinants of household participation is a central problem in financial economics. Why do certain households allocate capital to mutual funds and equities, while others restrict their savings to traditional, low-yield instruments like fixed deposits? The literature points to a combination of background risk, financial literacy, and trust in institutions. However, most existing studies rely on older data that predate the post-pandemic retail boom, the widespread adoption of digital KYC, and the systemic rise of social media financial influencers.

This paper addresses this gap by utilizing the comprehensive 2025 SEBI Household Survey. We pose four core research questions: (1) Who is likely to participate in the securities market? (2) What structural and behavioral barriers keep the majority of households out? (3) Which securities do participating households prefer, and for how long? and (4) Can we predict these outcomes given basic demographic and behavioral metrics?

The remainder of the paper is organized as follows: Section II details the dataset and methodology. Section III presents the empirical findings regarding determinants, behavioral anomalies, and securities preferences. Section IV details our predictive modeling framework. Section V synthesizes the findings, and Section VI translates them into targeted policy recommendations.

\section{Data and Methodology}

\subsection{Dataset}
We utilize micro-data from the 2025 SEBI Household Survey on Securities Market Participation, a nationally representative cross-sectional survey of $N = 109{,}430$ households administered across four geographic zones: North, South, East, and West India. Each observation corresponds to a unique household respondent, weighted for population representativeness via a pre-computed survey weight variable.

\subsection{Data Normalization and Schema Architecture}

\begin{figure}[htbp]
\centerline{\includegraphics[width=\linewidth]{../figures/data_layer.png}}
\caption{Architecture of the data layer showing the aggregation and normalization process of the respondent microdata.}
\label{fig:data_layer}
\end{figure}

To facilitate robust multivariate analysis, the raw survey microdata (originally comprising a single flat file with 449 columns) underwent a rigorous Extraction, Transformation, and Loading (ETL) process. Categorical survey responses and qualitative inputs were systematically transformed into strictly analytical data types (booleans, integers, and ordinal scales), eliminating unstructured free-text to ensure computational efficiency and statistical reliability.

As illustrated in Fig.~\ref{fig:data_layer}, the transformed data was subsequently loaded into a fully normalized PostgreSQL relational database. The schema architecture is anchored by a central \texttt{respondents} table, which establishes a unique primary key (\texttt{respondent\_id}) for each household and stores the population-representative survey weights. 

The remaining data is segmented across 16 interconnected, domain-specific tables encompassing geographic distribution, socioeconomic profiles, financial allocations, product awareness, and behavioral perceptions. Maintaining a strict one-to-one relationship with the core identity table, this modular and highly indexed design allows for rapid, optimized querying of specific behavioral intersections (such as isolating active equity investors to correlate objective financial literacy with risk tolerance), thereby forming the computational backbone for both our logistic regression and gradient boosting frameworks.

The sample is bifurcated into \textit{active investors} (those holding at least one financial security) ($n = 27{,}882$; 25.5\%) and \textit{non-investors} ($n = 81{,}548$; 74.5\%). Table~\ref{tab:descriptive} reports key demographic characteristics. The sample skews male (64.5\%) and is concentrated in younger cohorts: Gen~Z (18--27 years) and Millennials (28--43 years) together comprise 86.1\% of respondents. Nuclear family structures predominate (84.2\% of households).

\begin{table}[htbp]
\caption{Sample Composition and Participation Rates}
\label{tab:descriptive}
\centering
\begin{tabular}{lrrr}
\toprule
& \textbf{N} & \textbf{\% of Sample} & \textbf{Participation Rate} \\
\midrule
\multicolumn{4}{l}{\textit{Gender}} \\
\quad Male   & 70,616 & 64.5 & 30.8\% \\
\quad Female & 38,814 & 35.5 & 15.8\% \\
\midrule
\multicolumn{4}{l}{\textit{Life Stage}} \\
\quad Gen Z (18--27)       & 49,035 & 44.8 & 24.3\% \\
\quad Millennial (28--43)  & 45,235 & 41.3 & 28.1\% \\
\quad Gen X (44--59)       & 12,806 & 11.7 & 22.9\% \\
\quad Baby Boomer (60+)    &  2,354 &  2.2 & 13.3\% \\
\midrule
\multicolumn{4}{l}{\textit{Residence}} \\
\quad Urban & 76,931 & 70.3 & 28.8\% \\
\quad Rural & 32,499 & 29.7 & 17.6\% \\
\midrule
\multicolumn{4}{l}{\textit{Zone}} \\
\quad West  & 34,153 & 31.2 & 27.1\% \\
\quad South & 27,885 & 25.5 & 26.8\% \\
\quad North & 26,439 & 24.2 & 24.1\% \\
\quad East  & 20,953 & 19.1 & 22.8\% \\
\bottomrule
\end{tabular}
\end{table}

The survey employs a conditional routing design. Investors were directed to dedicated sub-modules (SS\_B3, SS\_B10) covering information sources and investment motivations, yielding $n = 9{,}797$ complete investor sub-responses. Non-investors were routed to a separate module (SS\_BB2) eliciting up to three self-reported barriers to participation, yielding $n = 14{,}794$ barrier responses. All group-level percentages for these modules are computed against the respective routing-eligible population, not the full sample.

\subsection{Exploratory Analysis}
The raw participation patterns in Table~\ref{tab:descriptive} reveal a market sharply stratified by socioeconomic position. The gender gap is the most visually immediate: male respondents participate at 30.8\% against 15.8\% for female respondents. The income gradient, however, is the steepest univariate signal: participation rises from 10.7\% in the lowest quintile (below INR~12,500/month) to 47.6\% in the highest (above INR~35,000/month), a more than fourfold difference. Urban households (28.8\%) participate at nearly double the rate of rural ones (17.6\%). Notably, Millennials (28--43 years) record the \textit{highest} participation rate of any cohort (28.1\%), exceeding Gen~X (22.9\%) and Gen~Z (24.3\%), suggesting the post-2020 retail boom was disproportionately captured by this age group rather than younger entrants.

\begin{figure}[htbp]
\centerline{\includegraphics[width=\linewidth]{../figures/fig_eda1_participation_demo.pdf}}
\caption{Market participation rates by gender, residence, and monthly income quintile. The income gradient (10.7\% to 47.6\%) is the steepest univariate signal, while the gender gap (30.8\% vs 15.8\%) persists across all income levels.}
\label{fig:participation_demo}
\end{figure}

A critical precondition for participation is product awareness. The survey's awareness module reveals a steep gradient: fixed deposits are known to 99.1\% of respondents, while mutual funds and ETFs reach 67.0\%, direct equity 62.0\%, and derivatives just 17.9\%. This awareness gap, between near-universal familiarity with traditional instruments and niche awareness of market-linked ones, maps closely onto the observed holdings distribution. Among the 27,882 active investors, MF/ETFs lead penetration (66.5\%), followed by FD/RD (57.2\%) and direct equity (49.3\%). The co-existence of MFs and FDs as the two most-held instruments indicates that Indian retail investors are not homogenously risk-seeking; the portfolio evidence points instead to a layered allocation model where market-linked products supplement rather than replace traditional savings.

\begin{figure}[htbp]
\centerline{\includegraphics[width=\linewidth]{../figures/fig_eda2_holdings.pdf}}
\caption{Instrument holdings penetration among active investors ($n=27{,}882$). MF/ETFs and Fixed Deposits co-dominate, indicating a layered allocation model rather than a shift away from traditional savings.}
\label{fig:holdings_penetration}
\end{figure}

The financial literacy data deepens this picture. When presented with nine objective financial knowledge questions, investors perform unevenly: while 80\%+ correctly identify that KYC can be completed online and that higher returns entail higher risk, a majority respond incorrectly or ``Don't Know'' to questions about compounding, diversification, and even basic SEBI tools (CAS, BSDA). This asymmetry (high procedural literacy, weak conceptual literacy) is directly relevant to the Dunning-Kruger analysis in Section~\ref{sec:results} and to the barriers analysis, where ``Lack of Knowledge'' ranks as the second most-cited deterrent among non-investors.

\subsection{Variable Operationalization}
The dependent variable $Y_i \in \{0, 1\}$ indicates whether household $i$ holds any financial security. We include twelve regressors across four groups. The \textit{awareness} block contains \textit{Product Awareness} ($n\_products\_aware$; count of financial instruments the respondent has heard of, 0--15, drawn from the Q21A multi-select module administered to all 109,430 respondents). The \textit{socioeconomic} block contains \textit{Education} (continuous years of formal schooling) and \textit{Income} (natural log of self-reported monthly household income in INR). The \textit{demographic} block contains \textit{Urban} (binary; urban vs.\ rural residence), \textit{Gender} (binary; male $= 1$), and \textit{Risk Tolerance} (ordinal; 1 = low, 2 = medium, 3 = high). The \textit{life stage} block uses Gen~Z as the reference category, with binary indicators for Millennial (28--43 years), Gen~X (44--59 years), and Baby Boomer (60+ years). The \textit{geographic} block uses the East zone as reference, with binary indicators for North, South, and West. All continuous regressors are standardized ($\mu = 0$, $\sigma = 1$) prior to estimation so that OR magnitudes are directly comparable across predictors.

\subsection{Econometric Specification}
We estimate a binary logistic regression via maximum likelihood:

\begin{equation}
\Pr(Y_i = 1) = \frac{1}{1 + e^{-(\beta_0 + \boldsymbol{\beta}^\top \mathbf{x}_i)}}
\label{eq:logit}
\end{equation}

\noindent where $\mathbf{x}_i$ is the vector of standardized regressors. Coefficients are exponentiated to yield odds ratios (ORs), interpreted as the multiplicative change in participation odds per one-standard-deviation increase in each predictor. Statistical significance is evaluated via Wald $z$-tests with heteroskedasticity-consistent (HC3) standard errors; survey weights are incorporated during estimation. We note that with $N \approx 100{,}000$, statistical significance is trivially achieved for any real effect; $p$-values are reported for completeness but the economically meaningful metrics are OR magnitude and confidence interval width. We include 12 regressors spanning awareness, socioeconomic, demographic, and geographic dimensions; two (Gen~X life stage and West zone) return ORs statistically indistinguishable from 1.0, providing a falsifiability check on the model.

\subsection{Predictive Framework}
To operationalize the regression findings, we implement Histogram-based Gradient Boosting Classifiers for three distinct prediction tasks: (Model~1) the binary likelihood of market participation; (Model~2) whether an investor holds market-linked instruments (equity or mutual funds) vs.\ traditional instruments only; and (Model~3) the propensity to hold equity over a long-term horizon. All models are trained on an 80/20 stratified split and evaluated via the area under the ROC curve (AUC).

\section{Empirical Results}
\subsection{Socioeconomic Determinants of Participation}\label{sec:results}
Table~\ref{tab:logit} presents the full logistic regression results ($N = 97{,}626$; McFadden pseudo-$R^2 = 0.290$, excellent fit range 0.20--0.40). Ten of twelve predictors are statistically significant; Gen~X (vs.\ Gen~Z) and West zone (vs.\ East) are not, with ORs of 0.993 and 1.021 respectively, confirming the model does not return spurious significance at this sample size.

\begin{table}[htbp]
\caption{Logistic Regression: Determinants of Market Participation}
\label{tab:logit}
\centering
\begin{tabular}{lcccc}
\toprule
\textbf{Predictor} & \textbf{OR} & \textbf{95\% CI} & \textbf{$z$} & \textbf{Sig.} \\
\midrule
\multicolumn{5}{l}{\textit{Awareness}} \\
\quad Product Awareness   & 2.218 & [2.166,\ 2.272] & 65.2 & *** \\
\midrule
\multicolumn{5}{l}{\textit{Socioeconomic}} \\
\quad Log(Income)         & 1.603 & [1.564,\ 1.643] & 37.4 & *** \\
\quad Education (yrs)     & 1.556 & [1.518,\ 1.594] & 35.6 & *** \\
\midrule
\multicolumn{5}{l}{\textit{Demographic}} \\
\quad Male                & 1.526 & [1.490,\ 1.564] & 33.9 & *** \\
\quad Urban               & 1.455 & [1.422,\ 1.489] & 32.0 & *** \\
\quad Risk Tolerance      & 1.170 & [1.145,\ 1.196] & 14.0 & *** \\
\midrule
\multicolumn{5}{l}{\textit{Life Stage (ref: Gen Z)}} \\
\quad Millennial          & 1.124 & [1.098,\ 1.151] &  9.8 & *** \\
\quad Gen X               & 0.993 & [0.969,\ 1.017] & $-$0.6 & \phantom{***} \\
\quad Baby Boomer         & 0.907 & [0.881,\ 0.934] & $-$6.6 & *** \\
\midrule
\multicolumn{5}{l}{\textit{Zone (ref: East)}} \\
\quad North               & 0.850 & [0.825,\ 0.876] & $-$10.7 & *** \\
\quad South               & 0.842 & [0.816,\ 0.869] & $-$10.7 & *** \\
\quad West                & 1.021 & [0.989,\ 1.054] &   1.3 & \phantom{***} \\
\midrule
\multicolumn{5}{l}{\textit{N} = 97,626 \quad McFadden pseudo-$R^2$ = 0.290 (excellent fit: 0.20--0.40)} \\
\multicolumn{5}{l}{*** $p < 0.001$. HC3 robust SEs. Survey-weighted GLM (Binomial).} \\
\bottomrule
\end{tabular}
\end{table}

As depicted in Fig.~\ref{fig:odds_ratios}, product awareness (OR = 2.22) is the dominant driver of participation, substantially exceeding income (1.60), education (1.56), male gender (1.53), and urban residency (1.45). The magnitude gap is meaningful: moving one standard deviation up in awareness predicts a larger change in participation odds than any structural or demographic variable in the model.

The partial mediation structure is evident in the OR magnitudes: income (1.60) and education (1.56) retain large independent effects even after awareness is controlled, yet the relatively modest gap between them and awareness (2.22) indicates that part of the socioeconomic advantage operates through knowing more products. Awareness is therefore not merely a socioeconomic proxy; it is a distinct, policy-actionable lever that educated, higher-income households happen to possess more of.

The life stage results are instructive: Millennials participate \textit{more} than Gen~Z (OR = 1.12), plausibly reflecting the post-2020 retail boom; Gen~X is statistically indistinguishable from Gen~Z; and Baby Boomers participate significantly less (OR = 0.91). On geography, both South (OR = 0.84) and North (OR = 0.85) zones underperform East after controls, despite South's higher raw participation rates, indicating that South's raw advantage is fully explained by its demographic composition. West is no different from East after controls.

\begin{figure}[htbp]
\centerline{\includegraphics[width=\linewidth]{../figures/fig1_odds_ratios.pdf}}
\caption{Forest plot of odds ratios for determinants of market participation. Baseline is 1.0. Product awareness ($n\_products\_aware$, OR = 2.22) is the single strongest predictor, followed by log-income (1.60) and education (1.56). All predictors standardised to unit-SD units.}
\label{fig:odds_ratios}
\end{figure}

\subsection{Barriers to Entry: Latent Archetype Analysis}
To better understand the 74.5\% of households that remain completely outside the securities market, we move beyond simple frequency rankings and cluster the 18 self-reported barriers into latent archetypes. Using K-means clustering on the subset of non-investors who identified specific deterrents ($n=14,705$), we identify three distinct profiles: \textit{Trust-Deficient} (46.3\%), \textit{Knowledge-Gated} (29.0\%), and \textit{Fear-Driven} (24.8\%).

The demographic cross-tabulation (Fig.~\ref{fig:archetype_demographics}) reveals that the \textit{Knowledge-Gated} profile is disproportionately concentrated in rural areas and among the lowest income quartile (Q1). In contrast, \textit{Trust-Deficiency} remains a persistent barrier that scales with wealth, remaining high even in urban higher-income cohorts. This segmentation implies that financial inclusion cannot be solved with a one-size-fits-all awareness campaign; it requires targeted interventions: educational for the rural knowledge-gated, and institutional/regulatory for the urban trust-deficient.


\begin{figure}[htbp]
\centerline{\includegraphics[width=\linewidth]{../figures/fig11_archetype_demographics.pdf}}
\caption{Demographic distribution of barrier archetypes.}
\label{fig:archetype_demographics}
\end{figure}

\subsection{Behavioral Anomalies: Overconfidence \& Dunning-Kruger}
Diving deeper into the ``Knowledge'' barrier, we analyzed the divergence between self-reported market familiarity and actual scores on an objective financial literacy quiz. We discovered a profound Dunning-Kruger effect: investors who are ``overconfident'' (high perceived familiarity, low actual knowledge) are disproportionately likely to hold high-risk assets like Futures \& Options (F\&O) and Cryptocurrency compared to calibrated investors (Fig. \ref{fig:dunning_kruger}). This behavioral anomaly underscores why traditional models of rational participation often fail to capture modern retail dynamics.

\begin{figure}[htbp]
\centerline{\includegraphics[width=\linewidth]{../figures/fig5_dunning_kruger.pdf}}
\caption{Risky Asset Penetration. Overconfident investors participate in highly speculative assets at significantly higher rates than calibrated investors.}
\label{fig:dunning_kruger}
\end{figure}

\subsection{Investor Education Programme (IEP) Exposure}
Among the 27,882 active investors, 2,768 (9.9\% of all investors) cite SEBI or industry Investor Education Programmes as one of their top information sources.

As shown in Fig.~\ref{fig:iep_analysis}, IEP-exposed investors score markedly higher on the nine-item objective financial literacy quiz relative to the broader non-IEP investor population (5.30 vs.\ 3.75, a 41\% gap). Product awareness is near-identical (10.23 vs.\ 10.23), indicating the knowledge advantage is specific to conceptual depth rather than product breadth. A before-after design would be needed to confirm causality, but the cross-sectional evidence is consistent with IEP deepening financial comprehension among participants.

The divergence in investment philosophy is equally pronounced: IEP-exposed investors are more likely to cite financial goals as a primary motive (30.2\% vs.\ 22.8\%) and less likely to seek quick gains (21.7\% vs.\ 23.3\%), suggesting that IEP reorients investors toward patient, goal-based capital formation. They also hold fewer Fixed Deposits (49.3\% vs.\ 55.1\%), consistent with a modest tilt toward market-linked instruments. Taken together, these findings position IEP as both a knowledge deepening and a philosophical reorientation intervention.

\begin{figure}[htbp]
\centerline{\includegraphics[width=\linewidth]{../figures/fig12_iep_analysis.pdf}}
\caption{IEP-exposed vs.\ non-IEP investors. IEP users score 41\% higher on financial literacy; they also show stronger goal orientation and lower quick-gains motivation.}
\label{fig:iep_analysis}
\end{figure}

\subsection{Securities Preferences, Demographics, and Influencers}
Once households overcome the barrier to entry, their product preferences shift dynamically with income. Fig. \ref{fig:securities_income} illustrates the penetration of various financial instruments across income quartiles. Fixed Deposits (FDs) remain the bedrock of Indian savings, showing uniformly high penetration. However, the adoption of market-linked securities (specifically Mutual Funds/ETFs and Direct Equity) scales dramatically in the higher income brackets, indicating that risk appetite scales non-linearly with wealth.

\begin{figure}[htbp]
\centerline{\includegraphics[width=\linewidth]{../figures/fig3_securities_income.pdf}}
\caption{Securities penetration among investors, segmented by income quartile. Market-linked instruments scale sharply with wealth.}
\label{fig:securities_income}
\end{figure}

Crucially, the choice of securities is heavily modulated by where investors source their information. Contrasting investors who rely exclusively on Social Media ``Finfluencers'' against those using traditional Financial Professionals reveals a sharp behavioral divide. As seen in Fig. \ref{fig:finfluencers}, the Finfluencer cohort is primarily motivated by ``Quick Gains,'' whereas the professionally-advised cohort seeks ``Long Term Growth.'' (Prevalence here is calculated as the percentage of investors within a specific information-source cohort who selected a given motive). This ties directly to the overconfidence trap discussed earlier, painting a picture of a segmented market where information source dictates risk behavior.

\begin{figure}[htbp]
\centerline{\includegraphics[width=\linewidth]{../figures/fig6_finfluencers.pdf}}
\caption{Investment Motives segmented by Information Source. Prevalence denotes the percentage of investors within each cohort holding a specific motive. Social media heavily biases investors toward short-term speculative gains.}
\label{fig:finfluencers}
\end{figure}

\subsection{Goal--Portfolio Consistency}

Beyond demographics and information sources, we test whether household portfolios are internally consistent with their stated primary financial goals. As shown in Fig.~\ref{fig:goal_alignment_heatmap}, households citing long-term objectives such as Retirement (72.4\%) and Children’s Education (77.2\%) exhibit strong alignment, with high penetration in mutual funds and provident fund instruments (PPF/EPF), indicating coherent long-horizon financial planning.

To formalize this, we define a \textit{goal--portfolio consistency} metric as the proportion of households whose primary goal maps to a theoretically aligned instrument set (e.g., Retirement $\rightarrow$ MF/EPF; Daily Trading $\rightarrow$ Equity/F\&O) and who hold at least one such instrument.

However, we identify a stark ``Say--Do'' paradox among speculative entrants. While households citing Daily Trading as their primary goal should, in principle, exhibit high exposure to equities and derivatives, we observe an internal consistency score of just 7.3\% (Fig.~\ref{fig:consistency_score}). This is not merely a mismatch but a structural phenomenon: households self-identify with a trading-oriented financial identity without participating in the asset classes that operationalize such behavior.

We therefore classify households into three segments: (i) \textit{aligned investors} (goals and portfolios match), (ii) \textit{constrained investors} (risk-oriented goals but conservative portfolios), and (iii) \textit{aspirational investors} (speculative self-identification without corresponding asset adoption). The ``daily trading'' cohort overwhelmingly falls into the third category.

This misalignment provides a structural explanation for the overconfidence patterns documented earlier. Aspirational investors exhibit high perceived market familiarity but lack both instrument exposure and technical literacy, creating the precise conditions for Dunning--Kruger effects. In practice, this segment is likely to engage with markets episodically and without adequate risk management, making it particularly vulnerable to adverse financial outcomes.

\begin{figure}[htbp]
\centerline{\includegraphics[width=\linewidth]{../figures/fig8_goal_alignment_heatmap.pdf}}
\caption{Asset penetration by primary financial goal. While long-term planners show high adoption of MFs and FDs, the speculative cohort exhibits surprisingly low adoption of their target instruments.}
\label{fig:goal_alignment_heatmap}
\end{figure}

\begin{figure}[htbp]
\centerline{\includegraphics[width=\linewidth]{../figures/fig9_consistency_score.pdf}}
\caption{Internal Consistency Score across primary goals. Consistency is measured as the presence of goal-appropriate instruments in the portfolio. The speculative 'Daily Trading' cohort shows a massive intent--action gap.}
\label{fig:consistency_score}
\end{figure}

\subsection{Investment Horizons}
Holding duration preferences (Fig. \ref{fig:duration}) challenge the notion that all retail investors are short-term speculators. A significant plurality of equity and mutual fund investors define their holding periods as ``Long term''. In contrast, sophisticated derivative products (F\&O) are predominantly utilized for ``Short term'' horizons, aligning with the motivations of the Finfluencer-driven cohort.

\begin{figure}[htbp]
\centerline{\includegraphics[width=\linewidth]{../figures/fig4_duration.pdf}}
\caption{Preferred holding durations across different asset classes. Direct equities are largely held for the long term, whereas derivatives skew short term.}
\label{fig:duration}
\end{figure}

\section{Predictive Modeling: Forecasting Household Participation}
To operationalize our findings, we utilized Histogram-based Gradient Boosting Classifiers to predict household behavior. All three models incorporate product awareness ($n\_products\_aware$) alongside standard demographic and socioeconomic features. For investor-specific models (Models 2 and 3), the feature space is further expanded to include the behavioral traits uncovered in our earlier analysis: perceived familiarity, objective financial literacy, and information source reliance.

\begin{figure*}[t]
\centerline{\includegraphics[width=0.9\textwidth]{../figures/fig7_predictive_importance.pdf}}
\caption{Normalized feature importances across the three predictive models.}
\label{fig:predictive_importance}
\end{figure*}

As shown in Fig.~\ref{fig:predictive_importance}, the feature importance profiles differ sharply across models: awareness dominates participation prediction, while familiarity and objective knowledge take over once inside the market.

\subsection{Model 1: Predicting Participation}
\textbf{Objective:} Predict the binary outcome of whether a household is an active investor ($1$) or a non-investor ($0$). \\
\textbf{Data Fields Used:} Years of Education, Log-transformed Monthly Income, Urban Residency (Binary), Gender (Binary), Risk Tolerance Preference Score, and Product Awareness ($n\_products\_aware$). \\
\textbf{Algorithm \& Split:} HistGradientBoostingClassifier (200 max iterations, balanced class weights) trained on an 80/20 train-test split. \\
\textbf{Performance:} The model achieved an AUC of \textbf{0.814}. \\
\textbf{Insight:} Product awareness is the single most predictive feature, consistent with its dominant OR of 2.22 in the logistic regression. Awareness of more instruments is independently predictive beyond any demographic factor.

\subsection{Model 2: Predicting Securities Choice}
\textbf{Objective:} For households that \textit{are} investors, predict if they will hold market-linked securities (Equities or Mutual Funds = $1$) versus sticking exclusively to traditional instruments like Fixed Deposits ($0$). \\
\textbf{Data Fields Used:} All features from Model 1 (including Product Awareness), \textbf{plus} Self-Perceived Stock Market Familiarity, Actual Financial Literacy Score, Social Media Reliance (Binary), and Professional Advice Reliance (Binary). \\
\textbf{Algorithm \& Split:} HistGradientBoostingClassifier trained on an 80/20 train-test split. \\
\textbf{Performance:} The model achieved an AUC of 0.820. \\
\textbf{Insight:} Predicting \textit{how} individuals participate requires behavioral context. Once we include perceived familiarity and objective knowledge, the model's discriminatory power surges.

\subsection{Model 3: Predicting Investment Horizon}
\textbf{Objective:} For equity investors, predict if they will define their holding period as ``Long Term'' ($1$) versus Short/Mid-term ($0$). \\
\textbf{Data Fields Used:} Same extended demographic, awareness, and behavioral feature set used in Model 2. \\
\textbf{Algorithm \& Split:} HistGradientBoostingClassifier trained on an 80/20 train-test split. \\
\textbf{Performance:} The model achieved an AUC of 0.612. \\
\textbf{Insight:} Forecasting \textit{how long} an investor will hold equity remains the hardest prediction task, suggesting that horizon is shaped by situational factors outside our cross-sectional feature space.

\subsection{External Validation and Geographic Opportunity Mapping}
To validate the real-world utility of our predictive framework, we cross-referenced the results of Model 1 with aggregate market registry data from the \texttt{state\_layer} module. This process involved applying the respondent-level model to state-level demographic averages (income, literacy, urbanization) to calculate a ``Demographic Potential'' score for each of India's 35 states and union territories.

As shown in Fig.~\ref{fig:geographic_validation}, the model's predictions correlate significantly with actual market participation rates ($r=0.491$, $p<0.01$). More importantly, this cross-validation allows us to identify the ``Participation Gap'', defined as the delta between a state's demographic readiness and its actual market penetration. 

\begin{figure}[htbp]
\centerline{\includegraphics[width=0.85\linewidth]{../figures/model_validation_gap_analysis.pdf}}
\caption{Geographic Validation of Predictive Model. (A) Regression of model-predicted potential against actual market registry data. (B) The Participation Gap Analysis identifies states where demographics suggest higher participation potential than currently observed (Red) vs over-performing states (Blue).}
\label{fig:geographic_validation}
\end{figure}

Our analysis reveals several ``Under-Penetrated'' regions including \textbf{Lakshadweep, Puducherry, Mizoram, and Kerala}, where high literacy and income levels suggest a latent capacity for participation that remains uncaptured by current market infrastructure. Conversely, states like \textbf{Delhi, Gujarat, and Haryana} exhibit a negative gap, indicating that market enthusiasm in these clusters significantly exceeds what would be expected based on basic demographics alone, likely driven by mature brokerage networks and localized financial cultures.

\vspace{0.3cm}
These predictive models validate our core thesis: awareness and behavioral context, not demographics alone, determine participation and post-entry choices. The geographic validation further suggests that targeted institutional outreach in ``Under-Penetrated'' clusters could unlock significant new retail capital.

\section{Conclusion}
The 2025 SEBI Household Survey reveals a securities market with a three-layer structure. At the entry layer, participation is primarily gated by product awareness (OR = 2.22), exceeding even income (1.60) and education (1.56) as a predictor. This finding reframes the participation problem: the dominant barrier is not poverty but unfamiliarity with the instrument landscape. The partial mediation structure confirms that education and income partly work \textit{through} awareness, meaning financial outreach that expands product knowledge is a complementary lever to income growth rather than a substitute for it.

At the post-entry layer, demographics cease to predict behavior. Instead, information source becomes the dominant axis of segmentation. Investors relying on social media finfluencers concentrate in quick-gains motives and speculative instruments, while professionally-advised investors orient toward long-term growth. The Dunning-Kruger analysis reinforces this: overconfident investors, high on perceived familiarity but low on objective knowledge, disproportionately hold Futures \& Options and cryptocurrency, instruments whose risk profiles they demonstrably misunderstand.

At the goal-alignment layer, a structural disconnect emerges between investor intent and portfolio reality. The ``Daily Trading'' cohort shows a goal-portfolio consistency score of just 7.3\%, indicating that a large share of self-identified speculators have not yet acquired the instruments needed to operationalize their stated strategy. These aspirational investors represent the segment most exposed to adverse outcomes: they self-identify with high-risk financial identities without the literacy or portfolio depth to manage the associated risks.

Taken together, these findings establish that broadening Indian market participation requires parallel interventions at all three layers: expanding awareness to unlock entry, improving information quality to correct post-entry behavior, and closing the intent-action gap to protect the aspirational cohort.

\section{Policy Implications}

\subsection{Prioritise Product Awareness Over Generic Literacy}
Product awareness is the strongest individual predictor of participation (OR = 2.22), yet the awareness gradient is steep: while fixed deposits are known to 99.1\% of respondents, mutual funds reach only 67.0\%, direct equity 62.0\%, and derivatives 17.9\%. Generic financial literacy programs that teach abstract concepts without naming specific instruments miss this gap. SEBI and AMFI should design campaigns that systematically introduce named products (MFs, ETFs, SIPs) to non-investor households, particularly in North and South zones where participation underperforms East even after demographic controls.

\subsection{Segment Non-Investor Interventions by Barrier Archetype}
The 74.5\% of households outside the market are not a homogeneous group. Archetype clustering identifies three profiles with distinct policy needs: \textit{Trust-Deficient} (46.3\%), who require institutional accountability measures such as visible investor grievance redressal, SEBI's Investor Protection Fund promotion, and regulated broker transparency; \textit{Knowledge-Gated} (29.0\%), concentrated in rural areas and the lowest income quartile, who need geographically targeted financial literacy with vernacular delivery; and \textit{Fear-Driven} (24.8\%), who respond to loss-framing education and could be nudged through first-loss protection products or guaranteed-return onramps (e.g., SGB, RBI Floating Rate Bonds) that reduce perceived downside risk before transitioning to equity.

\subsection{Scale and Reform Investor Education Programmes}
IEP-exposed investors score 41\% higher on objective financial literacy than the broader non-IEP investor population (5.30 vs.\ 3.75), and show a measurable shift toward goal-based investment motives (30.2\% vs.\ 22.8\%) and away from quick-gains seeking (21.7\% vs.\ 23.3\%). Yet only 9.9\% of all investors cite IEP as a top information source. Scaling IEP reach through mandatory employer-side programs, vernacular digital delivery, and partnerships with Jan Dhan account infrastructure could replicate these outcomes at population scale. Critically, IEP content should emphasize conceptual depth (compounding, diversification, instrument risk profiles) rather than procedural knowledge (KYC, demat), where literacy is already high. The goal-alignment data provides direct evidence that this approach works: investors with long-term goals (Retirement, Children's Education) exhibit portfolio consistency scores of 72.4\% and 77.2\% respectively, confirming that goal-based financial framing translates reliably into coherent instrument selection.

\subsection{Close the Gender Participation Gap}
Female investors participate at 15.8\% versus 30.8\% for male investors, a gap that persists after controlling for income and education (OR = 1.53 for male gender). This is the single largest identifiable untapped pool of potential market participants. Policy responses should move beyond generic inclusion rhetoric toward specific mechanisms: women-only SHG-linked investment collectives, female-targeted IEP content addressing trust deficits and loss aversion, and incentive structures in platforms that lower the minimum investment threshold for first-time women investors.

\subsection{Regulate the Finfluencer Information Channel}
Social media finfluencers are the second most-used information source (5.0\% of all investors; on par with friends and family at 5.1\%), and the finfluencer cohort is disproportionately motivated by quick gains relative to the professionally-advised cohort. This is associated with a disproportionate quick-gains orientation that, combined with the overconfidence patterns documented earlier, is likely to amplify speculative behavior. SEBI's existing finfluencer registration framework should be enforced more aggressively, with mandatory disclosure of qualifications, prohibited performance claims, and platform-level liability for unregistered investment advice.

\subsection{Gate Derivatives Access on Objective Literacy}
The Dunning-Kruger analysis shows that overconfident investors (those with high perceived familiarity but low objective knowledge) hold Futures \& Options and crypto at significantly higher rates than calibrated investors. The goal-alignment data reinforces this: the daily trading cohort, whose strategy requires active derivatives exposure, has a portfolio consistency score of just 7.3\%, suggesting they aspire to speculative trading without the instruments or knowledge to do so safely. SEBI should consider replacing the current self-certification suitability test for derivatives access with a brief objective literacy assessment, analogous to the approach used for sophisticated investor classification in other jurisdictions. This would filter out the most behaviorally at-risk segment before they encounter leveraged instruments.

\section{Appendix}

\subsection{AI Stockbroker Sales Helper}
The final phase of this research involved the implementation of an AI-powered Research Assistant, the \textit{DSM Stockbroker Sales Helper}, designed to operationalize our empirical findings into actionable sales intelligence. The agent architecture is designed to bridge the gap between static academic results and real-world sales strategy.

\subsubsection{System Architecture}
The agent is built on a multi-layered agentic framework:
\begin{itemize}
    \item \textbf{Reasoning Engine}: Powered by the Gemini 1.5 Pro large language model, the agent utilizes a \textit{LangGraph}-based state machine. This allows for iterative reasoning loops where the agent can plan a query, execute a tool, observe the result, and refine its strategy before providing a final answer.
    \item \textbf{Dynamic Data Access}: The agent is natively integrated with our normalized PostgreSQL database ($N=109{,}430$). Using a specialized SQL Toolkit, it can generate and execute complex queries in real-time to find real-world lead examples or calculate specific geographic aggregates.
    \item \textbf{Operational Tools}: \textit{LeadScorer} (Gradient Boosting wrapper), \textit{GeographicAnalyst} (State-level gap analysis), and \textit{PitchOptimizer} (Behavioral what-if simulations).
\end{itemize}

\subsection{Strategic Sales Dashboard}
To host the AI agent and visualize our empirical findings, we developed a high-fidelity web application built with Next.js and React. The dashboard is designed with a "Premium-First" aesthetic to match the high-end nature of financial research.

\subsubsection{Design System \& Aesthetics}
The dashboard utilizes a curated design system focused on editorial clarity and atmospheric depth:
\begin{itemize}
    \item \textbf{Typography}: We utilize a high-contrast typographic pairing of \textit{Fraunces} (an editorial serif) for narrative synthesis and a strictly defined monospace font for data-heavy tool outputs and SQL indicators.
    \item \textbf{Theming}: A custom "Paper" theme was implemented using OKLCH color variables, providing a sophisticated dark-mode aesthetic with semantic tokens for different information layers (\texttt{paper-1} through \texttt{paper-3}).
    \item \textbf{Glassmorphism}: The interface employs semi-transparent overlays and blurred background effects to create a sense of depth and focus, common in modern premium fintech applications.
\end{itemize}

\subsubsection{Interactive Components}
We implemented several custom components to enhance the sales workflow:
\begin{itemize}
    \item \textbf{Real-Time Status Indicators}: Pulsing emerald status dots and "Text Shimmer" animations provide immediate visual feedback on the agent's background processing state.
    \item \textbf{Research Audit Trail}: A dedicated sidebar log that persists "Checkmarked" tool indicators, allowing brokers to verify the agent's data sources at a glance.
    \item \textbf{Adaptive Multi-Line Input}: An intelligent text area that supports both rapid execution and complex, multi-paragraph lead descriptions.
\end{itemize}

\subsection{Microdata Explorer}
To provide full transparency into the SEBI 2025 dataset, we implemented a robust Data Explorer module. This interface serves as an administrative and research portal for auditing the raw microdata that powers our predictive models.

\subsubsection{Technical Architecture}
The explorer is built on a high-performance server-side processing engine:
\begin{itemize}
    \item \textbf{Dynamic Schema Discovery}: The system utilizes SQLAlchemy's inspection engine to dynamically reflect the database schema, allowing it to serve any microdata table without manual configuration.
    \item \textbf{Server-Side Sorting \& Filtering}: To maintain performance over hundreds of thousands of records, we implemented global search and sorting directly in the SQL layer. This allows for complex filtering across all columns (utilizing \texttt{ILIKE} pattern matching) with sub-millisecond response times.
    \item \textbf{HeroUI v3 Modular Table}: The interface leverages the advanced modular architecture of HeroUI v3, utilizing \texttt{Table.ScrollContainer} for native horizontal scrolling and strict accessibility compliance via the \texttt{isRowHeader} identifier for record keys.
\end{itemize}

\subsubsection{Visual Polish}
The Data Explorer extends the dashboard's premium aesthetic:
\begin{itemize}
    \item \textbf{Monospace Precision}: All raw data is rendered in a strictly defined monospace font at small optical sizes (\texttt{11px}), optimizing for information density while maintaining an "Audit-First" look.
    \item \textbf{Interactive Metrics}: A high-fidelity record counter and column indicator provide real-time feedback on the scope of the current data subset.
    \item \textbf{Debounced Search}: The global search bar utilizes a debounced input pattern to prevent unnecessary server load during rapid typing, providing a smooth "Sync-on-Type" experience.
\end{itemize}

\begin{thebibliography}{00}
\bibitem{b1} SEBI, ``Household Survey on Securities Market Participation 2025,'' Securities and Exchange Board of India, 2025.
\bibitem{b2} Kruger, Justin \& Dunning, David. (2000). Unskilled and Unaware of It: How Difficulties in Recognizing One's Own Incompetence Lead to Inflated Self-Assessments. Journal of Personality and Social Psychology. 77. 1121-34. 10.1037//0022-3514.77.6.1121.
\end{thebibliography}

\end{document}